
% \tableofcontents

\section{Wprowadzenie}

Wszystkie dane są generowane przez symulator, nie pochodzą z rzeczywistych czujników. System dostępny jest przez przeglądarkę pod adresem \texttt{http://localhost:8080}.

\section{Podstawowe funkcje}

\begin{itemize}
    \item Uruchamianie symulacji pożaru na mapie
    \item Sterowanie brygadami i patrolami
    \item Otrzymywanie rekomendacji AI
    \item Analiza danych z czujników
\end{itemize}

\section{Interfejs}

\begin{itemize}
    \item Mapa: pokazuje sektory, brygady, patrole, czujniki
    \item Panel boczny: info o sektorze, rekomendacje, sterowanie
    \item Przycisk "Uruchom symulację", "Zatrzymaj", "Załaduj konfigurację"
\end{itemize}

\subsection{Panele boczne}

\subsubsection{Panel informacyjny sektora}

Wyświetla się po kliknięciu na sektor:

\begin{itemize}
    \item \textbf{ID sektora}
    \item \textbf{Typ sektora:} FOREST, GRASS, WATER, BUILDING
    \item \textbf{Stan ognia:} INACTIVE, ACTIVE, LOST
    \item \textbf{Poziom ognia:} 0-100
    \item \textbf{Poziom spalenia:} 0-100\%
    \item \textbf{Poziom gaszenia:} 0-100\%
    \item \textbf{Przypisane brygady:} Lista ID brygad
    \item \textbf{Dane z czujników:} Temperatura, wilgotność, wiatr, CO2, PM2.5
\end{itemize}

\subsubsection{Panel rekomendacji}

Wyświetla rekomendacje systemu:

\begin{itemize}
    \item Lista rekomendowanych akcji
    \item Dla każdej rekomendacji: ID brygady, ID sektora, priorytet
    \item Przycisk "Zaakceptuj" dla każdej rekomendacji
\end{itemize}

\subsubsection{Panel sterowania}

Zawiera przyciski do zarządzania symulacją:

\begin{itemize}
    \item \textbf{Uruchom symulację:} Start nowej symulacji
    \item \textbf{Zatrzymaj symulację:} Stop bieżącej symulacji
\end{itemize}

\subsection{Warstwy danych}

\subsubsection{Warstwa sektorów}

Wyświetla wszystkie sektory lasu z ich stanami. Można włączyć/wyłączyć przez menu warstw.

\subsubsection{Warstwa brygad}

Wyświetla pozycje i stany brygad pożarniczych. Można filtrować po stanie (dostępne, w drodze, gaszące).

\subsubsection{Warstwa patroli}

Wyświetla pozycje i stany patroli leśnych. Można filtrować po stanie.

\subsubsection{Warstwa czujników}

Wyświetla lokalizacje czujników. Kliknięcie na czujnik wyświetla jego dane.

\subsubsection{Warstwa rekomendacji}

Wyświetla wizualizację rekomendacji na mapie (strzałki od brygad do sektorów docelowych).

\section{Najczęstsze problemy}

\begin{itemize}
    \item Brak czujników: nie wszystkie konfiguracje je mają
    \item Brak brygad: poczekaj aż się zwolnią
    \item Brak reakcji: sprawdź czy symulacja jest uruchomiona, odśwież stronę
    \item Brak rekomendacji: poczekaj, sprawdź czy są aktywne pożary
    \item Mapa nie ładuje się: sprawdź konfigurację i połączenie z backendem
\end{itemize}

\section{Ograniczenia}

\begin{itemize}
    \item System nie posiada alertów w czasie rzeczywistym
    \item Wszystkie dane są generowane przez symulator
    \item Jednostki i pożary są wirtualne
    \item System może działać wolniej przy dużych konfiguracjach
    \item Rekomendacje AI są tylko sugestiami
\end{itemize}

\section{Porady}

\begin{itemize}
    \item Monitoruj wiele sektorów
    \item Używaj rekomendacji jako wskazówek
    \item Przypisuj brygady do różnych pożarów
    \item Wysyłaj patrole do sąsiednich sektorów
\end{itemize}

\section{Podsumowanie}

System jest symulacją do nauki i testowania strategii. Dane są generowane przez symulator, a rekomendacje AI są tylko sugestiami. Użytkownik decyduje o działaniach.

\section{Konfiguracja zaawansowana (LLM)}

\subsection{Przegląd}

System wspiera integrację z Large Language Model (LLM) do generowania rekomendacji i podejmowania decyzji przez agentów. Ta funkcjonalność jest \textbf{opcjonalna} i wymaga konfiguracji.

\subsection{Jak włączyć LLM}

W module \texttt{fire-support} znajdź plik \texttt{.local.env} i zmień następujące linie:

\begin{lstlisting}[language=bash, basicstyle=\small]
# Zamiast:
#LLM_ENABLED=true
#ENABLE_LLM_COORDINATION=true
#ENABLE_AGENT_COMMUNICATION=true

# Ustaw:
LLM_ENABLED=true
ENABLE_LLM_COORDINATION=true
ENABLE_AGENT_COMMUNICATION=true
\end{lstlisting}

Wybierz model OpenAI (domyślnie: \texttt{gpt-4o-mini}):

\begin{lstlisting}[language=bash, basicstyle=\small]
OPENAI_MODEL=gpt-4o-mini
# Alternatywy: gpt-4, gpt-4-turbo
\end{lstlisting}

Wybierz tryb generowania rekomendacji i decyzji agentów:

\begin{lstlisting}[language=bash, basicstyle=\small]
RECOMMENDATION_MODE=llm      # llm, mcts, hybrid
AGENT_DECISION_MODE=llm      # llm, heuristic, hybrid
\end{lstlisting}

Uruchom ponownie moduł \texttt{fire-support}:

\begin{lstlisting}[language=bash, basicstyle=\small]
docker compose restart fire-support-service
\end{lstlisting}

\subsection{Przykładowe konfiguracje}

\subsubsection{Konfiguracja 1: Bez LLM (domyślna)}

\begin{lstlisting}[language=bash, basicstyle=\small]
RECOMMENDATION_MODE=mcts
AGENT_DECISION_MODE=heuristic
LLM_ENABLED=false
\end{lstlisting}

\textbf{Efekt:} System używa tylko algorytmu MCTS (szybko, < 100ms).

\subsubsection{Konfiguracja 2: Pełna integracja LLM}

\begin{lstlisting}[language=bash, basicstyle=\small]
OPENAI_MODEL=gpt-4o-mini
RECOMMENDATION_MODE=llm
AGENT_DECISION_MODE=llm
LLM_ENABLED=true
ENABLE_LLM_COORDINATION=true
ENABLE_AGENT_COMMUNICATION=true
\end{lstlisting}

\textbf{Efekt:} Wszystko kontrolowane przez LLM (wolniejsze, 1-3s na decyzję).

\subsubsection{Problem 1: Błąd API OpenAI}

\textbf{Objawy:}
\begin{itemize}
    \item Rekomendacje nie pojawiają się
    \item Błąd w logach: \texttt{OpenAI API error}
\end{itemize}

\textbf{Przyczyna:} Brak klucza API lub nieprawidłowy klucz.

\textbf{Rozwiązanie:}
\begin{enumerate}
    \item Sprawdź czy masz klucz API OpenAI
    \item Dodaj klucz do pliku \texttt{.local.env}: \texttt{OPENAI\_API\_KEY=sk-...}
    \item Restart serwisu: \texttt{docker compose restart fire-support-service}
\end{enumerate}

\subsubsection{Problem 2: LLM jest wolny}

\textbf{Objawy:}
\begin{itemize}
    \item Rekomendacje pojawiają się po 3-5 sekundach
    \item System wydaje się ``zawieszać''
\end{itemize}

\textbf{Przyczyna:} API OpenAI jest wolne lub model jest zbyt duży.

\textbf{Rozwiązanie:}
\begin{enumerate}
    \item Użyj mniejszego modelu: \texttt{OPENAI\_MODEL=gpt-4o-mini}
\end{enumerate}

\subsubsection{Problem 3: Rekomendacje LLM są dziwne}

\textbf{Objawy:}
\begin{itemize}
    \item LLM sugeruje nielogiczne akcje
    \item Brygady są wysyłane w złe miejsca
\end{itemize}

\textbf{Przyczyna:} LLM nie rozumie kontekstu symulacji lub model jest za mały.

\textbf{Rozwiązanie:}
\begin{enumerate}
    \item Użyj większego modelu: \texttt{OPENAI\_MODEL=gpt-4}
    \item Przełącz na tryb hybrydowy: \texttt{RECOMMENDATION\_MODE=hybrid}
    \item Wyłącz LLM i użyj MCTS: \texttt{LLM\_ENABLED=false}
\end{enumerate}

\textbf{Gdy coś nie działa:}

\begin{enumerate}
    \item \textbf{Sprawdź logi:} \texttt{logs/fire-support/application.log}
    \item \textbf{Sprawdź status serwisu:} \texttt{docker compose ps}
\end{enumerate}

\subsubsection{Dlaczego brygady nie reagują na komendy?}

\textbf{Przyczyny:}
\begin{itemize}
    \item Symulacja jest zatrzymana
    \item Błąd komunikacji z backendem
\end{itemize}

\subsection{Gdzie szukać logów}

\textbf{Logi systemowe:}
\begin{itemize}
    \item \texttt{logs/fire-backend/application.log}
    \item \texttt{logs/fire-simulation/application.log}
    \item \texttt{logs/fire-support/application.log}
    \item \texttt{logs/fire-configuration/application.log}
\end{itemize}

\textbf{Logi Docker:}
\begin{lstlisting}[language=bash, basicstyle=\small]
docker compose logs fire-backend-service
docker compose logs fire-simulation-service
docker compose logs fire-support-service
\end{lstlisting}

\textbf{Logi przeglądarki:}
\begin{enumerate}
    \item Otwórz konsolę (F12)
    \item Zakładka ``Console''
    \item Zakładka ``Network'' (sprawdź błędy HTTP)
\end{enumerate}

